% 文件开头添加条件判断
\newif\ifstandalone
\standalonetrue  % 默认独立编译模式

% 检测是否在主文档中
\ifdefined\maindoc
  \standalonefalse
\fi

\ifstandalone
  \documentclass[12pt,a4paper]{ctexart}

  \usepackage[utf8]{inputenc}
  \usepackage{ctex}
  \usepackage{amsmath,amssymb}
  \usepackage{graphicx}
  \usepackage{hyperref}
  \usepackage{geometry}
  \usepackage{bm}

  \geometry{left=2.5cm,right=2.5cm,top=2.5cm,bottom=2.5cm}

  \title{第八章\quad 机器人的神经网络控制}
  \author{}
  \date{}
\fi

\ifstandalone
  \begin{document}
  \maketitle
\fi

本章给出了使用神经网络进行控制器设计的框架。这些结构具有通用逼近性质，使其能够在未知系统的反馈控制中使用，而无需系统参数线性化的要求或寻找回归矩阵。本章给出了能够保证闭环稳定性和有界权重的反馈控制拓扑结构和权重调整算法。

\section{引言}

近年来，在设计模仿生物系统功能的反馈控制系统方面进行了大量工作~\cite{white1992,miller1991}。最近人们对"通用无模型控制器"产生了极大兴趣，这种控制器不需要被控对象的数学模型，而是模仿生物过程的功能，在线学习它们正在控制的系统，从而自动提高性能。相关技术包括模糊逻辑控制——它模仿语言和推理功能，以及人工神经网络——它基于相互连接节点的生物神经元结构。神经网络（NN）在分类、模式识别以及数字信号处理和其他领域的开环应用方面取得了巨大成功。严格的分析已经展示了如何选择神经网络拓扑结构和权重，例如，用于区分指定的示例模式。到目前为止，神经网络在开环应用中的理论和应用已经得到了很好的理解，因此神经网络已成为信号处理人员和计算机科学家工具库中的重要工具。

关于使用神经网络对未知对象进行反馈控制，已有大量文献。直到20世纪90年代，设计和分析技术都是临时性的，没有可重复的设计算法或稳定性和保证性能的证明。尽管如此，文献中出现的仿真结果仍显示了良好的性能。大多数早期方法使用标准的反向传播权重调整~\cite{werbos1992}，因为当时适用于闭环控制目的的严格调整算法推导尚不可用。许多神经网络设计技术模仿自适应控制方法，而自适应控制方法已有严格的分析结果可用~\cite{astrom1989,landau1979,goodwin1984}，提出了基于间接辨识的控制、逆动力学控制、串并联技术等神经网络反馈控制拓扑结构。

按照Maxwell、Lyapunov、A.N.Whitehead、von Bertalanffy等人从事控制系统理论研究的哲学，要解决这些问题，必须从被控系统的已有知识开始。Narendra~\cite{narendra1991,narendra1991b}和其他人~\cite{miller1991,white1992,werbos1992}通过研究闭环系统中神经网络的动力学行为，包括计算反向传播调整所需的梯度，开创了严格的神经网络控制应用。其他为闭环神经网络控制器提供严格分析和设计技术的团队包括：\cite{sanner1991}展示了如何在反馈控制中使用径向基函数神经网络；\cite{chen1992,chen1994}提供了基于死区方法的神经网络调整算法；\cite{polycarpou1991,polycarpou1996}使用了投影方法；\cite{lewis1993,lewis1995}使用了e-mod方法；\cite{sadegh1993}为离散时间系统提供了神经网络控制器；\cite{rovithakis1994}使用动态神经网络进行反馈控制。

本章给出了一种基于几篇博士论文和一系列已发表工作的神经网络控制器设计和分析方法，包括~\cite{lewis1999,kim1998}。这里讨论的控制结构是多环控制器，其中一些环路包含神经网络，外部是跟踪单位增益PD环路。有可重复的设计算法和系统性能保证，包括小跟踪误差和有界神经网络权重。研究表明，随着对被控系统不确定性的增加或性能要求的提高，神经网络控制器需要越来越多的结构。更深入的阐述见于~\cite{lewis1999}，其中还展示了与模糊逻辑系统的关系。

\section{神经网络背景}

本节提供了反馈控制所需的神经网络结构背景知识。更多细节请参见~\cite{haykin1994,lewis1999,lewis1998b}。使神经网络对反馈控制有用的两个关键特征是其通用逼近性质和学习能力，后者源于其权重是可调参数，可以更新以提高控制器性能。通用逼近性质是使非线性网络结构比自适应机器人控制器更适合机器人控制的主要特征，后者通常依赖于确定回归矩阵，而这又要求可调参数线性化（LIP）。

\subsection{多层神经网络}

一个多层神经网络如图8.2.1所示。该神经网络有两层可调权重，这里称为"两层"网络。该神经网络没有内部反馈连接，因此称为前馈网络；没有内部动态，因此称为静态网络。神经网络输出$\bm{y}$是一个具有$m$个分量的向量，由回忆方程根据具有$n$个分量的输入向量$\bm{x}$确定：
\begin{equation}\label{eq:8.2.1}
y_i = \sum_{j=1}^{L} w_{ij} \sigma\left(\sum_{k=1}^{n} v_{jk} x_k + \bar{v}_j\right) + \bar{w}_i, \quad i = 1, \ldots, m
\end{equation}
其中$\sigma(\cdot)$是激活函数，$L$是隐层神经元数量。第一层互联权重用$v_{jk}$表示，第二层互联权重用$w_{ij}$表示。阈值偏移用$\bar{v}_j$、$\bar{w}_i$表示。

有许多不同的激活函数$\sigma(\cdot)$在广泛使用。对于使用多层神经网络的反馈控制，要求$\sigma(\cdot)$足够光滑，使得至少其一阶导数存在。合适的选择包括sigmoid函数
\begin{equation}\label{eq:8.2.2}
\sigma(z) = \frac{1}{1 + e^{-z}}
\end{equation}
双曲正切函数
\begin{equation}\label{eq:8.2.3}
\sigma(z) = \frac{e^z - e^{-z}}{e^z + e^{-z}}
\end{equation}
以及其他逻辑曲线型函数，还有高斯函数和其他各种函数。

通过将所有神经网络权重$v_{jk}$、$w_{ij}$收集到权重矩阵$\bm{V}^T$、$\bm{W}^T$中，神经网络回忆方程可以用向量形式写成：
\begin{equation}\label{eq:8.2.4}
\bm{y} = \bm{W}^T \sigma(\bm{V}^T \bm{x})
\end{equation}
阈值包含在权重矩阵$\bm{W}^T$、$\bm{V}^T$的第一列中；为了适应这一点，向量$\bm{x}$和$\sigma(\cdot)$需要通过在前面放置一个"1"来增广（例如$\bm{x} = [1, x_1, x_2, \ldots, x_n]^T$）。在这个方程中，为了表示式~\eqref{eq:8.2.1}，如果$\sigma(\cdot)$取为从$\mathbb{R}^L$到$\mathbb{R}^L$的对角函数，即对于向量$\bm{z} = [z_1, z_2, \ldots, z_L]^T \in \mathbb{R}^L$有$\sigma(\bm{z}) = \text{diag}\{\sigma(z_j)\}$，则具有足够的通用性。

注意，回忆方程在第一层权重和阈值$\bm{V}$方面是非线性的。

\textbf{通用函数逼近性质}。神经网络满足许多重要性质。对于反馈控制目的而言，主要关注的是通用函数逼近性质~\cite{hornik1989}。设$f(\bm{x})$是一个从$\mathbb{R}^n$到$\mathbb{R}^m$的一般光滑函数。那么可以证明，只要$\bm{x}$被限制在$\mathbb{R}^n$的紧集$S$内，就存在权重和阈值，使得
\begin{equation}\label{eq:8.2.5}
f(\bm{x}) = \bm{W}^T \sigma(\bm{V}^T \bm{x}) + \varepsilon
\end{equation}
对于一定数量的隐层神经元$L$。这对于一大类激活函数都成立，包括刚才提到的那些。值$\varepsilon$称为神经网络函数逼近误差，它通常随着网络规模$L$的增加而减小。事实上，对于任何选择的正数$\varepsilon_M$，可以找到前馈神经网络使得
\begin{equation}\label{eq:8.2.6}
\|\varepsilon\| < \varepsilon_M
\end{equation}
对于$S$中所有$\bm{x}$。为给定精度$\varepsilon_M$选择网络规模$L$的问题对于一般无结构全连接神经网络仍是一个开放问题。对于CMAC、FL网络和其他结构化非线性网络，这个问题可以随后解决。

最佳逼近给定非线性函数$f(\bm{x})$的理想神经网络权重矩阵$\bm{W}$、$\bm{V}$难以确定。事实上，它们甚至可能不唯一。然而，对于控制目的，人们只需要知道，对于指定的$\varepsilon_M$值，存在某些理想逼近神经网络权重。然后，可以用下式给出$f(\bm{x})$的估计
\begin{equation}\label{eq:8.2.7}
\hat{f}(\bm{x}) = \hat{\bm{W}}^T \sigma(\hat{\bm{V}}^T \bm{x})
\end{equation}
其中$\hat{\bm{W}}$、$\hat{\bm{V}}$是理想神经网络权重的估计，由随后详细介绍的在线权重调整算法提供。注意，所有"带帽"量都是已知的，因为它们可以用可测量信号和当前神经网络权重估计来计算。

存在使逼近性质成立的理想权重的假设非常类似于自适应控制中的各种类似假设~\cite{astrom1989,landau1979}，包括Erzberger假设和参数线性化。非常重要的区别在于，在神经网络情况下，通用逼近性质总是成立的，而在自适应控制中，这些假设在实践中往往不成立，因此它们意味着对被控系统形式的限制。

\begin{quote}
\textit{注：读者可以通过}\href{https://frankjie09.github.io/Robot-Manipulator-Control-CN/figures/Neural_Network_Control.html}{\textbf{交互式可视化工具}}\textit{观察神经网络学习非线性函数的过程，理解隐藏层神经元数量、学习率对逼近精度的影响。}
\end{quote}

\textbf{克服NLIP问题}。多层神经网络在权重$\bm{V}$方面是非线性的，因此产生保证闭环反馈系统中稳定性和有界权重的权重调整算法直到几年前才被发现。如果采取正确和严格的方法，NLIP问题很容易克服。一种方法是以下对泰勒级数的适当使用~\cite{lewis1996}。

定义函数估计误差、权重估计误差和隐层输出误差
\begin{align}
\tilde{f} &= f - \hat{f} = f - \hat{\bm{W}}^T \sigma(\hat{\bm{V}}^T \bm{x}) \label{eq:8.2.8} \\
\tilde{\bm{W}} &= \bm{W} - \hat{\bm{W}}, \quad \tilde{\bm{V}} = \bm{V} - \hat{\bm{V}} \label{eq:8.2.9} \\
\tilde{\sigma} &= \sigma - \hat{\sigma} = \sigma(\bm{V}^T \bm{x}) - \sigma(\hat{\bm{V}}^T \bm{x}) \label{eq:8.2.10}
\end{align}

对于任何$\bm{z}$，可以写出泰勒级数
\begin{align}
\sigma(\bm{z}) &= \sigma(\hat{\bm{z}}) + \sigma'(\hat{\bm{z}})\tilde{\bm{z}} + O(\tilde{\bm{z}})^2 \label{eq:8.2.11} \\
\intertext{其中$\sigma'$是雅可比矩阵，最后一项表示$\tilde{\bm{z}}^2$阶的项。因此}
\tilde{\sigma} &= \sigma'(\hat{\bm{V}}^T \bm{x})\tilde{\bm{V}}^T \bm{x} + O(\tilde{\bm{V}}^T \bm{x})^2 \label{eq:8.2.12}
\end{align}

这个关键方程允许将函数估计误差写成
\begin{equation}\label{eq:8.2.13}
\tilde{f} = \tilde{\bm{W}}^T \left[\hat{\sigma} - \hat{\sigma}'\hat{\bm{V}}^T \bm{x}\right] + \hat{\bm{W}}^T \hat{\sigma}' \tilde{\bm{V}}^T \bm{x} + w
\end{equation}
其中$w$包含高阶项。第一项有$\tilde{\bm{W}}$乘以已知量（方括号内），第二项有$\tilde{\bm{V}}$乘以已知量。当随后在闭环误差动力学中使用时，这种形式允许导出保证闭环稳定性的$\bm{V}$和$\bm{W}$的调整律。

\subsection{线性参数神经网络}

如果式~\eqref{eq:8.2.4}中的第一层权重和阈值$\bm{V}$固定，只调整第二层权重和阈值$\bm{W}$，那么神经网络只有一层可调权重。可以定义固定函数$\phi(\bm{x}) = \sigma(\bm{V}^T \bm{x})$，使得这种单层神经网络具有回忆方程
\begin{equation}\label{eq:8.2.14}
y = \bm{W}^T \phi(\bm{x})
\end{equation}
其中$\phi(\cdot): \mathbb{R}^n \to \mathbb{R}^L$，$L$是隐层神经元数量。该神经网络在可调参数$\bm{W}$方面是线性的，因此更容易调整。事实上，可以使用标准的自适应控制证明来导出合适的调整算法。

虽然是LIP，但这些神经网络相比要求确定回归矩阵的标准自适应控制方法仍有巨大优势，因为如果正确选择$\phi(\cdot)$函数，它们满足通用逼近性质。注意，自适应控制器在系统参数方面是线性的，而单层神经网络在神经网络权重方面是线性的。单层神经网络相比双层神经网络的一个优势是，对于前者存在关于为指定逼近精度选择隐层神经元数量$L$的确定结果。LIP网络的一个缺点是Barron~\cite{barron1993}已经证明了逼近精度的下界。

\textbf{功能连接基神经网络}。如果$\sigma(\cdot)$不是对角矩阵，而是允许$\phi(\cdot)$是从$\mathbb{R}^n$到$\mathbb{R}^L$的一般函数，则获得更大的通用性。这称为功能连接神经网络（FLNN）~\cite{sadegh1993}。对于LIP神经网络，函数逼近性质通常不成立。然而，只要将激活函数$\phi(\cdot)$选择为基，单层神经网络仍然可以逼近函数，基必须在$\mathbb{R}^n$的紧单连通集$S$上满足以下两个要求：

\begin{enumerate}
    \item $S$上的常数函数可以表示为式~\eqref{eq:8.2.14}的形式，对于有限数量$L$的隐层神经元。
    \item 式~\eqref{eq:8.2.14}的函数值域在从$S$到$\mathbb{R}^m$的连续函数空间中是稠密的，对于可数$L$。
\end{enumerate}

如果$\phi(\cdot)$提供基，那么从$\mathbb{R}^n$到$\mathbb{R}^m$的光滑函数$f(\bm{x})$可以在$\mathbb{R}^n$的紧集$S$上逼近为
\begin{equation}\label{eq:8.2.15}
f(\bm{x}) = \bm{W}^T \phi(\bm{x}) + \varepsilon
\end{equation}
对于某些理想权重和阈值$\bm{W}$和某些隐层神经元数量$L$。事实上，对于任何选择的正数$\varepsilon_M$，可以找到前馈神经网络使得$\|\varepsilon\| < \varepsilon_M$对于$S$中所有$\bm{x}$。

$f(\bm{x})$的逼近由下式给出
\begin{equation}\label{eq:8.2.16}
\hat{f}(\bm{x}) = \hat{\bm{W}}^T \phi(\bm{x})
\end{equation}
其中$\hat{\bm{W}}$是由调整算法给出的神经网络权重和阈值估计。对于LIP神经网络，函数估计误差由下式给出
\begin{equation}\label{eq:8.2.17}
\tilde{f} = \tilde{\bm{W}}^T \phi(\bm{x}) + \varepsilon
\end{equation}
当随后在误差动力学中使用时，这直接产生了保证闭环稳定性的$\bm{W}$的调整算法。

Barron~\cite{barron1993}已经证明，对于所有LIP逼近器存在一个基本下界，因此$\varepsilon$被$1/L^{2/n}$阶的项下界所限制。因此，随着神经网络输入$n$的增加，增加$L$以提高逼近精度变得不那么有效。如下所示，这对于使用单层神经网络设计的自适应控制器不是一个主要限制。这个下界问题在NLIP多层非线性网络中不存在。

我们经常用$\sigma(\cdot)$代替$\phi(\cdot)$，理解对于LIP网络，这个激活函数向量不是对角矩阵，而是从$\mathbb{R}^L$到$\mathbb{R}^L$的一般函数。

\textbf{高斯或径向基函数（RBF）网络}。在各种结构化非线性网络中，包括径向基函数、CMAC和模糊逻辑网络，激活函数基集的选择大大简化。这里将展示这种结构化非线性网络设计的关键在于比标准方程~\eqref{eq:8.2.1}中允许的更一般的神经网络阈值集，以及它们的适当选择。

常用的神经网络激活函数是高斯或径向基函数（RBF）~\cite{sanner1991}，当$\bm{x}$为标量时由下式给出
\begin{equation}\label{eq:8.2.18}
\sigma(x) = \exp\left(-\frac{(x-\bar{x})^2}{2p}\right)
\end{equation}
其中$\bar{x}$是均值，$p$是方差。RBF神经网络可以写成式~\eqref{eq:8.2.4}的形式，但相比通常的sigmoid神经网络有一个优势，即$n$维高斯函数从概率论、卡尔曼滤波和其他地方得到了很好的理解，因此$n$维RBF易于概念化。

第$j$个激活函数可以写成
\begin{equation}\label{eq:8.2.19}
\sigma_j(\bm{x}) = \exp\left[-(\bm{x}-\bar{\bm{x}}_j)^T \bm{P}_j^{-1}(\bm{x}-\bar{\bm{x}}_j)\right]
\end{equation}
其中$\bm{P}_j = \text{diag}\{p_{jk}\}$。定义激活函数向量为$\sigma(\bm{x}) \equiv [\sigma_1(\bm{x}), \sigma_2(\bm{x}), \ldots, \sigma_L(\bm{x})]^T$。如果协方差矩阵是对角的，使得$\bm{P}_j = \text{diag}\{p_{jk}\}$，那么式~\eqref{eq:8.2.19}变为可分离的，可以分解为分量形式
\begin{equation}\label{eq:8.2.20}
\sigma_j(\bm{x}) = \prod_{k=1}^{n} \exp\left[-\frac{(x_k-\bar{x}_{jk})^2}{2p_{jk}}\right] = \prod_{k=1}^{n} \sigma_{jk}(x_k)
\end{equation}
其中$\sigma_{jk}(x_k)$是$\sigma_j(\bm{x})$的第$k$个分量。因此，$n$维激活函数是$n$个标量函数的乘积。注意，这个方程具有式~\eqref{eq:8.2.1}中激活函数的形式，但有更一般的阈值，因为在每个隐层神经元$j$处，$\bm{x}$的每个不同分量都需要一个阈值；即图8.2.1中每个隐层神经元的阈值是一个向量。RBF方差$p_{jk}$和偏移$\bar{x}_{jk}$通常在RBF神经网络设计中选择并固定；通常只调整输出层权重$\bm{W}^T$。因此，RBF神经网络是一种特殊的FLNN（式~\eqref{eq:8.2.14}）（其中$\phi(\bm{x}) = \sigma(\bm{x})$）。

图8.2.2显示了$\bm{x} \in \mathbb{R}^2$情况下的可分离高斯函数。在这个图中，所有方差$p_{jk}$相同，均值$\bar{x}_{jk}$以特殊方式选择，在2维网格的节点点上间隔激活函数。为了形成在区域$\{-1 < x_1 \leq 1, -1 < x_2 \leq 1\}$上逼近函数的RBF神经网络，这里选择$L = 5 \times 5 = 25$个隐层神经元，沿$x_1$有5个单元，沿$x_2$有5个单元。其中九个神经元具有2维高斯激活函数，而沿边界的需要图示的"单侧"激活函数。

\section{使用静态神经网络的跟踪控制}

本节讨论使用非线性网络和假设全状态变量反馈的反馈跟踪控制设计。更多细节请参见~\cite{lewis1999}和其他列出的参考文献。如果全状态反馈可用，则静态前馈神经网络足以进行控制。这里提供了保证闭环稳定性和有界权重的控制拓扑和网络调整算法。所讨论的技术适用于一般非线性网络，包括神经网络和模糊逻辑系统~\cite{wang1997,lewis1997}，因此缩写NN此后可理解为指"非线性网络"。所得的多环控制拓扑具有外部PD跟踪环路和内部神经网络反馈线性化或动作生成环路。发现反向传播调整通常不够，需要修改的调整算法来保证闭环性能。

许多工业机械系统以及汽车、飞机和航天器具有拉格朗日形式的动力学，其典型代表是刚性机器人系统。因此，将考虑拉格朗日机器人动力学~\cite{lewis1993,lewis1996}。这里介绍的神经网络控制技术也可应用于其他未知系统，包括某些重要的非线性系统类别~\cite{lewis1999}。

\subsection{机器人手臂动力学和误差系统}

刚性拉格朗日系统的动力学，包括机器人手臂，具有一些重要的物理、结构和无源性性质~\cite{craig1988,lewis1993,slotine1988,spong1989}，使得在其控制中使用神经网络非常自然。在设计任何控制器时都应考虑这些性质——事实上，它们为神经网络控制器的严格设计算法提供了基础。

$n$连杆刚性（即无柔性连杆或高频关节/电机动力学）机器人操作臂的动力学可以用拉格朗日形式表示为
\begin{equation}\label{eq:8.3.1}
M(\bm{q})\ddot{\bm{q}} + V_m(\bm{q},\dot{\bm{q}})\dot{\bm{q}} + G(\bm{q}) + F(\dot{\bm{q}}) + \bm{\tau}_d = \bm{\tau}
\end{equation}
其中$\bm{q}(t) \in \mathbb{R}^n$是关节变量向量，其元素是机器人手臂关节角度或连杆伸长。$M(\bm{q})$是惯性矩阵，$V_m(\bm{q},\dot{\bm{q}})$是科里奥利/向心矩阵，$G(\bm{q})$是重力向量，$F(\dot{\bm{q}})$是摩擦力。有界的未知扰动（包括例如非结构化未建模动力学）用$\bm{\tau}_d$表示，控制输入力矩是$\bm{\tau}(t)$。机器人动力学具有以下标准性质：

\textbf{性质1}：$M(\bm{q})$是正定对称矩阵，有界于$m_1\bm{I} \leq M(\bm{q}) \leq m_2\bm{I}$，其中$m_1$、$m_2$是正常数。

\textbf{性质2}：矩阵$V_m(\bm{q},\dot{\bm{q}})$的范数被$v_b(\|\bm{q}\|)\|\dot{\bm{q}}\|$所限制，对于某个函数$v_b(\bm{q})$。

\textbf{性质3}：矩阵$\dot{M} - 2V_m$是斜对称的。这等价于内部力不做功的事实。

\textbf{性质4}：未知扰动满足$\|\bm{\tau}_d\| \leq d_B$，其中$d_B$是正常数。

给定期望的手臂轨迹$\bm{q}_d(t) \in \mathbb{R}^n$，跟踪误差为
\begin{equation}\label{eq:8.3.2}
\bm{e}(t) = \bm{q}_d(t) - \bm{q}(t)
\end{equation}
滤波跟踪误差为
\begin{equation}\label{eq:8.3.3}
\bm{r}(t) = \dot{\bm{e}}(t) + \bm{\Lambda}\bm{e}(t)
\end{equation}
其中$\bm{\Lambda}$是对称正定设计参数矩阵，通常选择为对角矩阵。如果找到控制器使得$\bm{r}(t)$有界，则$\bm{e}(t)$也有界；事实上$\|\bm{e}\| \leq \frac{\|\bm{r}\|}{\sigma_{\min}(\bm{\Lambda})}$，其中$\sigma_{\min}(\bm{\Lambda})$是$\bm{\Lambda}$的最小奇异值。

对$\bm{r}(t)$求导并使用式~\eqref{eq:8.3.1}，可以用滤波跟踪误差的形式写出手臂动力学
\begin{equation}\label{eq:8.3.4}
M\dot{\bm{r}} = -V_m\bm{r} - \bm{\tau} + \bm{f}(\bm{x})
\end{equation}
其中非线性机器人函数为
\begin{equation}\label{eq:8.3.5}
\bm{f}(\bm{x}) = M(\bm{q})\left(\ddot{\bm{q}}_d + \bm{\Lambda}\dot{\bm{e}}\right) + V_m(\bm{q},\dot{\bm{q}})\left(\dot{\bm{q}}_d + \bm{\Lambda}\bm{e}\right) + G(\bm{q}) + F(\dot{\bm{q}})
\end{equation}

计算$\bm{f}(\bm{x})$所需的向量$\bm{x}$可以定义为，例如
\begin{equation}\label{eq:8.3.6}
\bm{x} = \begin{bmatrix} \bm{e}^T & \dot{\bm{e}}^T & \bm{q}_d^T & \dot{\bm{q}}_d^T & \ddot{\bm{q}}_d^T \end{bmatrix}^T
\end{equation}
这是可测量的。函数$\bm{f}(\bm{x})$包含潜在未知的机器人参数，包括负载质量和复杂的摩擦形式。

用于轨迹跟踪的合适控制输入由计算力矩型控制给出
\begin{equation}\label{eq:8.3.7}
\bm{\tau} = \hat{\bm{f}}(\bm{x}) + K_v\bm{r} - \bm{v}(t)
\end{equation}
其中$K_v$是增益矩阵，通常选择为对角矩阵，$\hat{\bm{f}}(\bm{x})$是由某种方式提供的机器人函数$\bm{f}(\bm{x})$的估计。鲁棒化信号$\bm{v}(t)$需要补偿未建模非结构化扰动。使用这个控制，闭环误差动力学为
\begin{equation}\label{eq:8.3.8}
M\dot{\bm{r}} = -(V_m + K_v)\bm{r} + \tilde{\bm{f}} + \bm{\tau}_d + \bm{v}
\end{equation}

在计算控制信号时，估计$\hat{\bm{f}}$可以由几种技术提供，包括自适应控制或神经或模糊网络。辅助控制信号$\bm{v}(t)$可以由几种技术选择，包括滑模方法和其他鲁棒控制方法。

假设期望轨迹有界，使得
\begin{equation}\label{eq:8.3.9}
\begin{bmatrix} \bm{q}_d \\ \dot{\bm{q}}_d \\ \ddot{\bm{q}}_d \end{bmatrix} \leq \bm{q}_B
\end{equation}
其中$\bm{q}_B$是已知标量界。容易证明对于每个时刻$t$，$\bm{x}(t)$被限制为
\begin{equation}\label{eq:8.3.10}
\|\bm{x}\| \leq c_0 + c_1\|\bm{r}\| + c_2\|\bm{r}\|^2 \equiv q_B
\end{equation}
对于可计算的正常数$c_0$、$c_1$、$c_2$。

\subsection{自适应控制}

机器人学中的标准自适应控制技术~\cite{craig1988,slotine1988,spong1989,lewis1993}假设非线性机器人函数在可调参数方面是线性的，使得
\begin{equation}\label{eq:8.3.11}
\bm{f}(\bm{x}) = \bm{W}^T \phi(\bm{x}) + \varepsilon
\end{equation}
其中$\bm{W}$是未知参数向量，$\phi(\bm{x})$是已知回归矩阵。控制选择为
\begin{equation}\label{eq:8.3.12}
\bm{\tau} = \hat{\bm{W}}^T \phi(\bm{x}) + K_v\bm{r} - \bm{v}(t)
\end{equation}
并为参数估计$\hat{\bm{W}}$确定调整算法。这由LIP假设促进。如果逼近误差$\varepsilon$非零，则必须修改调整算法以保证有界的参数估计。可以使用各种鲁棒化技术，包括$\sigma$-mod~\cite{ioannou1984}、e-modification~\cite{narendra1987}或死区技术~\cite{kreisselmeier1986}。现在已经存在不需要LIP或确定回归矩阵的自适应控制技术~\cite{colbaugh1994}。将这些与这里将要讨论的神经网络控制器的复杂性进行比较是有趣的。

\subsection{神经网络反馈跟踪控制器}

神经网络将用于式~\eqref{eq:8.3.7}中的控制$\bm{\tau}(t)$，为未知机器人函数$\bm{f}(\bm{x})$提供估计$\hat{\bm{f}}$。神经网络逼近性质保证总是存在神经网络能在给定精度$\varepsilon_N$内完成这一点。对于2层NLIP神经网络，逼近为
\begin{equation}\label{eq:8.3.13}
\bm{f}(\bm{x}) = \bm{W}^T \sigma(\bm{V}^T \bm{x}) + \varepsilon
\end{equation}
而对于1层LIP神经网络，它是
\begin{equation}\label{eq:8.3.14}
\bm{f}(\bm{x}) = \bm{W}^T \phi(\bm{x}) + \varepsilon
\end{equation}
其中$\phi(\cdot)$选择为基。

神经网络控制器的结构出现在图8.3.1中，其中$\bm{e} \equiv [\bm{e}^T, \dot{\bm{e}}^T]^T$。为$\bm{f}(\bm{x})$提供估计的神经网络出现在内部控制环路中，外部有PD项$K_v\bm{r}$提供的跟踪环路。在控制理论术语中，内部环路是反馈线性化控制器~\cite{slotine1991}，而在计算机科学术语中，它是动作生成环路~\cite{werbos1992,miller1991,white1992}。这种多环智能控制结构自然地从机器人控制概念中导出，不是临时性的。因此，它不受关于合适神经网络控制拓扑的哲学讨论的影响，包括关于前馈与反馈、直接 versus 间接等的常见讨论。值得注意的是，图中的静态前馈神经网络通过在周围闭合反馈环路变成了动态神经网络（参见~\cite{narendra1991}）。

\textbf{神经网络相比LIP自适应控制的优势}。每个机器人系统都有其自己的回归矩阵，因此必须为每个系统确定不同的$\phi(\bm{x})$。这个回归矩阵通常很复杂，确定它可能很耗时。人们注意到，回归矩阵有效地为特定给定系统提供了函数逼近的基。另一方面，非线性网络的通用逼近性质表明，神经网络为所有足够光滑的系统提供了基。因此，神经网络可以用于逼近所有刚性机器人系统的光滑$\bm{f}(\bm{x})$，有效地允许为一类系统设计单个可调控制器。无需确定回归矩阵，因为相同的神经网络激活函数足以应对整个对象类。

即使是LIP的1层神经网络也有这个优势。注意，1层神经网络在系统参数方面不是线性的，而是在神经网络权重方面是线性的。即使是LIP神经网络也具有通用逼近性质，因此不需要回归矩阵。

\subsection{初始跟踪误差和初始神经网络权重}

现在需要确定如何调整神经网络权重以产生保证的闭环稳定性。本节考虑神经网络控制器设计的几种情况。所有情况都需要以下构造。

由于神经网络逼近性质在紧集上成立，必须如下定义允许的初始条件集。让神经网络逼近性质对于式~\eqref{eq:8.3.5}中给出的函数$\bm{f}(\bm{x})$在式~\eqref{eq:8.2.6}中具有给定精度$\varepsilon_N$，对于半径$\bm{b}_x > \bm{q}_B$的球内所有$\bm{x}$成立。定义允许的初始跟踪误差集为
\begin{equation}\label{eq:8.3.15}
S_r = \{\bm{r}(0): \|\bm{r}(0)\| \leq \frac{(b_x - q_B)}{c_0}\}
\end{equation}

注意，神经网络的逼近精度决定了$S_r$的大小。对于更大的神经网络（即更多隐层单元），$\varepsilon_N$在更大半径$\bm{b}_x$下较小。因此，允许的初始条件集$S_r$更大。另一方面，更活跃的期望轨迹（例如包含更高频率分量）导致更大的加速度$\ddot{\bm{q}}_d(t)$，产生更大的界$\bm{q}_B$，从而减小$S_r$。值得注意的是$S_r$依赖于PD设计比$\bm{\Lambda}$——$c_0$和$c_2$都依赖于$\bm{\Lambda}$。

我们初始条件要求的一个关键特征是其独立于神经网络初始权重。这与文献中其他技术形成鲜明对比，在那些技术中，稳定性证明依赖于选择某些初始稳定神经网络权重，这是非常困难的。

\section{线性参数神经网络的调参算法}

假设现在使用LIP FLNN根据式~\eqref{eq:8.2.15}逼近非线性机器人函数~\eqref{eq:8.3.5}，其中$\|\varepsilon\| \leq \varepsilon_M$在紧集上，理想逼近权重$\bm{W}$是常数，$\phi(\bm{x})$选择为基。$\bm{f}(\bm{x})$的估计由式~\eqref{eq:8.3.14}给出。那么控制律~\eqref{eq:8.3.7}变为
\begin{equation}\label{eq:8.4.1}
\bm{\tau} = \hat{\bm{W}}^T \phi(\bm{x}) + K_v\bm{r} - \bm{v}(t)
\end{equation}
闭环滤波误差动力学~\eqref{eq:8.3.8}为
\begin{equation}\label{eq:8.4.2}
M\dot{\bm{r}} = -(V_m + K_v)\bm{r} + \tilde{\bm{W}}^T\phi(\bm{x}) + \bm{\tau}_d + \varepsilon + \bm{v}
\end{equation}

下一个定理导出了用e-mod项增强的神经网络控制器的调整律，如~\cite{narendra1987}中所述。

现在必须假设理想权重$\bm{W}$是常数且有界，使得
\begin{equation}\label{eq:8.4.3}
\|\bm{W}\|_F \leq W_B
\end{equation}
其中界$W_B$已知。在~\cite{polycarpou1996}中，展示了如何使用标准自适应鲁棒技术来避免假设界$W_B$已知。

\textbf{定理8.4-1}：（神经网络权重调整算法）。

设期望轨迹$\bm{q}_d(t)$被$\bm{q}_B$所界，初始跟踪误差$\bm{r}(0)$在$S_r$中。设神经网络重构误差界$\varepsilon_N$和扰动界$d_B$是常数。假设理想神经网络目标权重被$W_B$所界，如式~\eqref{eq:8.4.3}。设机器人手臂的控制输入由式~\eqref{eq:8.4.1}给出，$\bm{v}(t)=0$，增益满足
\begin{equation}\label{eq:8.4.cond1}
K_{v\min} > \frac{(W_B + \frac{\varepsilon_N + d_B}{\|\phi\|})^2}{4}
\end{equation}

设权重调整为
\begin{equation}\label{eq:8.4.cond2}
\dot{\hat{\bm{W}}} = \bm{F}\phi(\bm{x})\bm{r}^T - \kappa\bm{F}\|\bm{r}\|\hat{\bm{W}}
\end{equation}
其中$\bm{F}=\bm{F}^T>0$，$\kappa>0$是小设计参数。那么滤波跟踪误差$\bm{r}(t)$和神经网络权重估计$\hat{\bm{W}}(t)$是一致最终有界（UUB）的，实际界分别由式~\eqref{eq:8.4.bound8}、\eqref{eq:8.4.bound9}的右侧给出。而且，通过增加跟踪增益$K_v$，跟踪误差可以变得任意小。

\textbf{证明}：

让神经网络逼近性质~\eqref{eq:8.2.15}对于式~\eqref{eq:8.3.5}中给出的函数$\bm{f}(\bm{x})$在紧集$\|\bm{x}\| \leq \bm{b}_x$上具有给定精度$\varepsilon_N$，其中$\bm{b}_x > \bm{q}_B$。设$\bm{r}(0) \in S_r$。那么在时刻0逼近性质成立。

定义Lyapunov函数候选并求导
\begin{align}
L &= \frac{1}{2}\bm{r}^T M(\bm{q})\bm{r} + \frac{1}{2}\text{tr}\{\tilde{\bm{W}}^T \bm{F}^{-1} \tilde{\bm{W}}\} \label{eq:8.4.lyap3} \\
\dot{L} &= \bm{r}^T M \dot{\bm{r}} + \frac{1}{2}\bm{r}^T \dot{M} \bm{r} + \text{tr}\{\tilde{\bm{W}}^T \bm{F}^{-1} \dot{\tilde{\bm{W}}}\} \label{eq:8.4.deriv4} \\
\intertext{代入式~\eqref{eq:8.4.2}得}
\dot{L} &= -\bm{r}^T (V_m + K_v)\bm{r} + \frac{1}{2}\bm{r}^T \dot{M} \bm{r} + \text{tr}\{\tilde{\bm{W}}^T(\bm{F}^{-1}\dot{\tilde{\bm{W}}} + \phi\bm{r}^T)\} + \bm{r}^T(\bm{\tau}_d + \varepsilon) \label{eq:8.4.subst5} \\
\intertext{然后，使用调整规则~\eqref{eq:8.4.cond2}得}
\dot{L} &= -\bm{r}^T K_v \bm{r} + \kappa \|\bm{r}\| \text{tr}\{\tilde{\bm{W}}^T(\bm{W} - \tilde{\bm{W}})\} + \bm{r}^T(\bm{\tau}_d + \varepsilon) \label{eq:8.4.tune6}
\end{align}

由于
\begin{align}
\text{tr}\{\tilde{\bm{W}}^T(\bm{W} - \tilde{\bm{W}})\} &= \langle \tilde{\bm{W}}, \bm{W} \rangle - \|\tilde{\bm{W}}\|_F^2 \\
&\leq \|\tilde{\bm{W}}\|_F \|\bm{W}\|_F - \|\tilde{\bm{W}}\|_F^2
\end{align}

得到
\begin{equation}\label{eq:8.4.result7}
\dot{L} \leq -K_{v\min}\|\bm{r}\|^2 + \kappa\|\bm{r}\|\|\tilde{\bm{W}}\|_F(W_B - \|\tilde{\bm{W}}\|_F) + \|\bm{r}\|(d_B + \varepsilon_M)
\end{equation}

这在花括号中的项为正时是负的。完成平方得
\begin{equation}
\dot{L} \leq -\|\bm{r}\|\left[K_{v\min}\|\bm{r}\| - \kappa W_B \|\tilde{\bm{W}}\|_F + \kappa\|\tilde{\bm{W}}\|_F^2 - (d_B + \varepsilon_M)\right]
\end{equation}

当以下条件时保证为正
\begin{equation}\label{eq:8.4.bound8}
\|\bm{r}\| \geq \frac{\kappa W_B^2/4 + (d_B + \varepsilon_M)}{K_{v\min}} \equiv b_r
\end{equation}
或
\begin{equation}\label{eq:8.4.bound9}
\|\tilde{\bm{W}}\|_F \geq W_B/2 + \sqrt{\kappa W_B^2/4 + (d_B + \varepsilon_M)/\kappa} \equiv b_W
\end{equation}

因此，$\dot{L}$在紧集外是负的。根据增益选择式~\eqref{eq:8.4.cond1}确保由$\|\bm{r}\| \leq b_r$定义的紧集包含在$S_r$中，因此逼近性质在整个过程中成立。这证明了$\|\bm{r}\|$和$\|\bm{W}\|_F$的UUB。

这个证明类似于~\cite{narendra1987}，但修改为反映神经网络只在紧集上逼近的事实。在这个结果中，现在需要理想权重$\bm{W}$的界$W_B$，它出现在PD增益条件~\eqref{eq:8.4.cond1}中。现在发现了跟踪误差$\bm{r}(t)$和权重$\bm{W}(t)$的界。注意，使用修改的权重调整算法建立$\bm{W}$的界不需要持续激励（PE）条件。

神经网络权重调整算法~\eqref{eq:8.4.cond2}由反向传播时间项加上一个新的第二项组成。$\kappa$项（称为e-modification）的重要性在于它在$\|\tilde{\bm{W}}\|_F$中添加了二次项，使得可以证明$\dot{L}$在$(\|\bm{r}\|, \|\tilde{\bm{W}}\|_F)$平面上的紧集外为负~\cite{narendra1987}。e-mod项使调整律对未建模动力学具有鲁棒性，因此不需要PE条件。在~\cite{polycarpou1991}中，使用投影算法来保持神经网络权重有界。在~\cite{chen1992,chen1994}中，采用了死区技术。

\textbf{权重初始化和在线调整}。本文导出的神经网络控制方案中没有预先的离线学习阶段。权重简单地初始化为零，因为那时图8.3.1显示控制器只是一个PD控制器。机器人学文献中的标准结果~\cite{dawson1990}表明，如果$K_v$足够大，PD控制器给出有界误差。因此，闭环系统保持稳定，直到神经网络开始学习。权重在系统跟踪期望轨迹时实时在线调整。随着神经网络学习$\bm{f}(\bm{x})$，跟踪性能提高。这比其他神经网络控制技术有重大改进，在那些技术中必须找到某些初始稳定权重，对于复杂非线性系统通常是不可能的。

\textbf{跟踪误差和神经网络权重估计误差的界}。式~\eqref{eq:8.4.bound8}的右侧可以视为跟踪误差的实际界，即$\bm{r}(t)$永远不会偏离它太远。重要的是从这个方程注意到，跟踪误差随着神经网络重构误差$\varepsilon_N$和机器人扰动$d_B$而增加，但通过选择大增益$K_v$，可以获得任意小的跟踪误差。尽管Barron对$\varepsilon$有下界限制，也是如此。

注意，神经网络权重$\hat{\bm{W}}$不能保证接近理想未知权重$\bm{W}$，后者给出$\bm{f}(\bm{x})$的良好逼近。然而，正如证明所保证的，只要$\hat{\bm{W}} - \bm{W}$有界，这就无关紧要。这保证了有界控制输入，从而可以实现跟踪目标。

\section{非线性参数神经网络的调参算法}

现在假设使用2层神经网络根据式~\eqref{eq:8.2.5}逼近机器人函数，其中$\|\varepsilon\| \leq \varepsilon_M$在紧集上，理想逼近权重$\bm{W}$、$\bm{V}$是常数。那么控制律~\eqref{eq:8.3.7}变为
\begin{equation}\label{eq:8.5.1}
\bm{\tau} = \hat{\bm{W}}^T \sigma(\hat{\bm{V}}^T \bm{x}) + K_v\bm{r} - \bm{v}(t)
\end{equation}

提出的神经网络控制结构如图8.3.1所示。

NLIP神经网络控制器比1层FLNN强大得多，因为不需要选择激活函数基。实际上，权重$\bm{V}$在线调整以自动提供合适的基。控制在第一层权重$\bm{V}$方面是非线性的，这在导出合适的权重调整律时带来复杂性。辅助信号$\bm{v}(t)$是随后详细介绍的函数，它提供面对NLIP问题产生的泰勒级数中高阶项的鲁棒性。

使用结果~\eqref{eq:8.2.13}正确处理NLIP问题，闭环误差动力学~\eqref{eq:8.3.8}可以写成
\begin{equation}\label{eq:8.5.2}
M\dot{\bm{r}} = -(V_m + K_v)\bm{r} + \tilde{\bm{W}}^T\hat{\sigma}'\hat{\bm{V}}^T\bm{x} + \tilde{\bm{W}}^T\hat{\sigma} + \bm{w} + \bm{\tau}_d + \bm{v}
\end{equation}
其中扰动项为
\begin{equation}\label{eq:8.5.3}
\bm{w} = \tilde{\bm{W}}^T\hat{\sigma}'\tilde{\bm{V}}^T\bm{x} + \bm{W}^T O(\tilde{\bm{V}}^T\bm{x})^2
\end{equation}

定义
\begin{align}
\bm{Z} &\equiv \begin{bmatrix} \bm{W} & 0 \\ 0 & \bm{V} \end{bmatrix}, \quad \hat{\bm{Z}} \equiv \begin{bmatrix} \hat{\bm{W}} & 0 \\ 0 & \hat{\bm{V}} \end{bmatrix}, \quad \tilde{\bm{Z}} \equiv \bm{Z} - \hat{\bm{Z}} \label{eq:8.5.4} \\
\intertext{和$\bm{Z}$、$\hat{\bm{Z}}$等价。不难证明}
\|\bm{w}\| &\leq C_0 + C_1\|\tilde{\bm{Z}}\|_F + C_2\|\tilde{\bm{Z}}\|_F\|\bm{r}\| \label{eq:8.5.5}
\end{align}
其中$C_i$是已知正常数。

假设理想权重有界，使得
\begin{equation}\label{eq:8.5.6}
\|\bm{Z}\|_F \leq Z_B
\end{equation}
其中$Z_B$已知。（可以使用标准自适应鲁棒控制技术来放宽$Z_B$已知的假设~\cite{polycarpou1996}。）

下一个定理给出了本章最强大的神经网络控制器。

\textbf{定理8.5-1}：（多层神经网络控制器的增强反向传播权重调整）。

设期望轨迹$\bm{q}_d(t)$被$\bm{q}_B$所界，初始跟踪误差$\bm{r}(0)$在$S_r$中。设理想目标神经网络权重被已知$Z_B$所界。取机器人动力学~\eqref{eq:8.3.1}的控制输入为式~\eqref{eq:8.5.1}，PD增益满足
\begin{equation}\label{eq:8.5.cond1}
K_{v\min} > \frac{(C_0 + C_1 Z_B + C_2 Z_B^2/4)(C_3 + Z_B)}{4}
\end{equation}
其中$C_3$在证明中定义。设鲁棒化项为
\begin{equation}\label{eq:8.5.cond2}
\bm{v}(t) = -K_z(\|\hat{\bm{Z}}\|_F + Z_B)\bm{r}
\end{equation}
增益为
\begin{equation}\label{eq:8.5.cond3}
K_z > C_2
\end{equation}

设神经网络权重调整为
\begin{equation}\label{eq:8.5.cond4}
\dot{\hat{\bm{W}}} = \bm{F}\hat{\sigma}'\hat{\bm{V}}^T\bm{x}\bm{r}^T - \bm{F}\hat{\sigma}'\tilde{\bm{V}}^T\bm{x}\bm{r}^T - \kappa\|\bm{r}\|\bm{F}\hat{\bm{W}}
\end{equation}
\begin{equation}\label{eq:8.5.cond5}
\dot{\hat{\bm{V}}} = \bm{G}\bm{x}(\hat{\bm{W}}^T\hat{\sigma}')^T\bm{r} - \kappa\|\bm{r}\|\bm{G}\hat{\bm{V}}
\end{equation}
其中任意常数矩阵$\bm{F}=\bm{F}^T>0$，$\bm{G}=\bm{G}^T>0$，$\kappa>0$是小标量设计参数。那么滤波跟踪误差$\bm{r}(t)$和神经网络权重估计$\hat{\bm{V}}$、$\hat{\bm{W}}$是UUB的，界分别由式~\eqref{eq:8.5.bound10}和\eqref{eq:8.5.bound11}具体给出。而且，通过增加式~\eqref{eq:8.5.1}中的增益$K_v$，跟踪误差可以保持任意小。

\textbf{证明}：

让神经网络逼近性质~\eqref{eq:8.2.5}对于式~\eqref{eq:8.3.5}中给出的函数$\bm{f}(\bm{x})$在紧集$\|\bm{x}\| \leq \bm{b}_x$上具有给定精度$\varepsilon_N$，其中$\bm{b}_x > \bm{q}_B$。设$\bm{r}(0) \in S_r$。那么在时刻0逼近性质成立。

定义Lyapunov函数候选并求导
\begin{align}
L &= \frac{1}{2}\bm{r}^T M(\bm{q})\bm{r} + \frac{1}{2}\text{tr}\{\tilde{\bm{W}}^T \bm{F}^{-1} \tilde{\bm{W}}\} + \frac{1}{2}\text{tr}\{\tilde{\bm{V}}^T \bm{G}^{-1} \tilde{\bm{V}}\} \label{eq:8.5.lyap6} \\
\dot{L} &= \bm{r}^T M \dot{\bm{r}} + \frac{1}{2}\bm{r}^T \dot{M} \bm{r} + \text{tr}\{\tilde{\bm{W}}^T \bm{F}^{-1} \dot{\tilde{\bm{W}}}\} + \text{tr}\{\tilde{\bm{V}}^T \bm{G}^{-1} \dot{\tilde{\bm{V}}}\} \label{eq:8.5.deriv7} \\
\intertext{从误差系统~\eqref{eq:8.5.2}代入得}
\begin{split}
\dot{L} &= -\bm{r}^T K_v \bm{r} + \frac{1}{2}\bm{r}^T (\dot{M} - 2V_m)\bm{r} - \bm{r}^T\bm{v} + \bm{r}^T(\bm{w} + \bm{\tau}_d) \\
&\quad + \text{tr}\{\tilde{\bm{W}}^T(\bm{F}^{-1}\dot{\hat{\bm{W}}} + \hat{\sigma}\bm{r}^T - \hat{\sigma}'\hat{\bm{V}}^T\bm{x}\bm{r}^T)\} \\
&\quad + \text{tr}\{\tilde{\bm{V}}^T(\bm{G}^{-1}\dot{\hat{\bm{V}}} + \bm{x}\bm{r}^T\hat{\bm{W}}^T\hat{\sigma}')\}
\end{split} \label{eq:8.5.subst8}
\end{align}

调整规则给出
\begin{equation}
\dot{L} = -\bm{r}^T K_v \bm{r} + \kappa\|\bm{r}\|\text{tr}\{\tilde{\bm{Z}}^T(\bm{Z}-\tilde{\bm{Z}})\} + \bm{r}^T(\bm{w}+\bm{\tau}_d+\varepsilon)
\end{equation}

由于
\begin{align}
\text{tr}\{\tilde{\bm{Z}}^T(\bm{Z}-\tilde{\bm{Z}})\} &= \langle \tilde{\bm{Z}}, \bm{Z} \rangle - \|\tilde{\bm{Z}}\|_F^2 \\
&\leq \|\tilde{\bm{Z}}\|_F \|\bm{Z}\|_F - \|\tilde{\bm{Z}}\|_F^2
\end{align}

得到
\begin{equation}
\dot{L} \leq -K_{v\min}\|\bm{r}\|^2 - K_z\|\bm{r}\|^2(\|\hat{\bm{Z}}\|_F + Z_B) + \kappa\|\bm{r}\|\|\tilde{\bm{Z}}\|_F(Z_B - \|\tilde{\bm{Z}}\|_F) + \|\bm{r}\|(C_0 + C_1\|\tilde{\bm{Z}}\|_F + C_2\|\tilde{\bm{Z}}\|_F\|\bm{r}\|)
\end{equation}

其中$K_{v\min}$是$K_v$的最小奇异值，最后一个不等式由于式~\eqref{eq:8.5.cond3}成立。

$\dot{L}$在花括号中的项为正时是负的。定义$C_3 = Z_B + C_1/\kappa$并完成平方得
\begin{equation}\label{eq:8.5.compact9}
\dot{L} \leq -\|\bm{r}\|\left[K_{v\min}\|\bm{r}\| - \kappa C_3^2/4 - C_0 - \frac{(C_3+Z_B/2)^2}{K_z-C_2}\right]
\end{equation}

当以下条件时保证为正
\begin{equation}\label{eq:8.5.bound10}
\|\bm{r}\| \geq \frac{\kappa C_3^2/4 + C_0}{K_{v\min}} + \frac{(C_3+Z_B/2)^2}{K_{v\min}(K_z-C_2)} \equiv b_r
\end{equation}

或
\begin{equation}\label{eq:8.5.bound11}
\|\tilde{\bm{Z}}\|_F \geq C_3/2 + \sqrt{C_3^2/4 + C_0/\kappa + (C_3+Z_B/2)^2/[\kappa(K_z-C_2)]} \equiv b_Z
\end{equation}

因此，$\dot{L}$在紧集外是负的。根据LaSalle扩展~\cite{lewis1993}，这证明了在该集合内控制保持有效时$\|\bm{r}\|$和$\|\bm{Z}\|_F$的UUB。然而，PD增益条件~\eqref{eq:8.5.cond1}表明由$\|\bm{r}\| \leq b_r$定义的紧集包含在$S_r$中，因此逼近性质在整个过程中成立。

调整算法~\eqref{eq:8.5.cond4}、~\eqref{eq:8.5.cond5}中的第一项不过是跟踪误差通过时间的反向传播。最后一项是Narendra e-mod，式~\eqref{eq:8.5.cond4}中的第二项是由于NLIP而产生的新型二阶项。鲁棒化项$\bm{v}(t)$也是由于NLIP问题而需要的。

定理8.4.1之后的评论在这里也是适用的。具体来说，没有预先离线调整阶段，神经网络权重调整在实时控制动作的同时在线发生。权重初始化在证明中不是问题，无需提供初始稳定权重；权重简单地初始化使得神经网络输出等于零，因为那时PD环路保持系统稳定，直到神经网络开始学习。然而，实际实验~\cite{gutierrez1998}表明，适当初始化权重很重要。一个好的选择是随机选择$\hat{\bm{V}}(0)$，$\hat{\bm{W}}(0)$等于零。

在实践中不需要计算常数$c_0$、$c_1$、$c_2$、$C_0$、$C_1$、$C_2$、$Z_B$或确定$S_r$。简单地选择大的神经网络规模$L$和PD增益$K_v$，然后执行仿真。然后用不同的$L$和$K_v$重复仿真。基于两次仿真之间行为的差异，可以修改$L$和$K_v$。

\textbf{直接反向传播权重调整}。反向传播算法（通常以离散时间形式）在文献中被无数次提出用于闭环控制。可以证明，如果神经网络逼近误差$\varepsilon$为零，扰动$\bm{\tau}_d(t)$为零，且泰勒级数~\eqref{eq:8.2.12}中没有高阶项，那么直接反向传播算法
\begin{equation}\label{eq:8.5.7}
\dot{\hat{\bm{W}}} = \bm{F}\hat{\sigma}'\hat{\bm{V}}^T\bm{x}\bm{r}^T, \quad \dot{\hat{\bm{V}}} = \bm{G}\bm{x}(\hat{\bm{W}}^T\hat{\sigma}')^T\bm{r}
\end{equation}
可以用于闭环控制代替式~\eqref{eq:8.5.cond4}、~\eqref{eq:8.5.cond5}。需要使用PE来证明使用直接反向传播时神经网络权重是有界的。提到的条件非常严格，在实践中不会发生；它们本质上要求机器人手臂是线性的（例如单连杆）。而且，在神经网络中不容易保证PE条件。

\subsection{神经网络控制器的无源性性质}

本文中使用的神经网络是静态前馈网络，但由于它们出现在闭环反馈中并使用微分方程调整，它们是动态系统。因此，可以讨论这些神经网络的无源性。一般来说，不能保证神经网络是被动的。然而，本文中的神经网络控制器具有一些重要的无源性性质，从而产生鲁棒的闭环性能。

系统$\dot{\bm{x}}=\bm{f}(\bm{x},\bm{u})$，$\bm{y}=\bm{h}(\bm{x},\bm{u})$如果验证所谓的功率形式等式~\cite{slotine1991}
\begin{equation}\label{eq:8.5.8}
\dot{L} = \bm{u}^T\bm{y} - g(\bm{x},\bm{u})
\end{equation}
对于某个下有界$\dot{L}(t)$和某个$g(t) \geq 0$，则称为被动的。即
\begin{equation}\label{eq:8.5.9}
\int_0^T \bm{u}^T\bm{y} \, dt \geq \int_0^T g(t) \, dt - \gamma_0^2
\end{equation}
对于所有$T \geq 0$和某个$\gamma_0^2 \geq 0$。如果系统是被动的，且此外
\begin{equation}\label{eq:8.5.10}
\int_0^T \bm{u}^T\bm{y} \, dt \geq \int_0^T g(t) \, dt - \gamma_0^2
\end{equation}
其中$g(t) > 0$，则系统称为耗散的。

如果$g(t)$是具有有界系数的$\|\bm{x}\|$的二次函数，则发生一种特殊的耗散性。我们称之为状态严格无源性（SSP）；那么，内部状态的范数以输送到系统的功率为界。令人惊讶的是，SSP概念在文献中没有得到广泛使用~\cite{lewis1993,seron1994}，尽管参见~\cite{goodwin1984}其中定义了输入和输出严格无源性。SSP结果在研究神经网络控制器的无源性性质中是关键的，并允许在没有持续激励假设的情况下得出某些内部有界性质。

\textbf{机器人跟踪误差动力学的无源性}

本文中的误差动力学（例如~\eqref{eq:8.5.2}）具有形式
\begin{equation}\label{eq:8.5.11}
M\dot{\bm{r}} = -(V_m + K_v)\bm{r} + \bm{\xi}_0
\end{equation}
其中$\bm{r}(t)$是滤波跟踪误差，$\bm{\xi}_0(t)$是适当定义的。该系统满足以下强无源性性质。

\textbf{定理8.5-2}：（机器人误差动力学的SSP）。

动力学~\eqref{eq:8.5.11}从$\bm{\xi}_0(t)$到$\bm{r}(t)$是状态严格被动系统。

\textbf{证明}：

取非负函数
\begin{equation}
L = \frac{1}{2}\bm{r}^T M(\bm{q})\bm{r}
\end{equation}

所以，使用~\eqref{eq:8.5.11}，得到
\begin{equation}
\dot{L} = -\bm{r}^T(V_m + K_v)\bm{r} + \bm{r}^T\bm{\xi}_0
\end{equation}

因此，斜对称性质产生功率形式
\begin{equation}
\dot{L} = \bm{r}^T\bm{\xi}_0 - \bm{r}^T K_v \bm{r}
\end{equation}

这是输送到系统的功率减去$\|\bm{r}\|^2$中的二次项，验证了状态严格无源性。

\textbf{2层神经网络控制器的无源性性质}

闭环系统如图8.3.1所示，其中神经网络出现在内部反馈线性化环路中。2层神经网络控制器的误差动力学由式~\eqref{eq:8.5.2}给出。闭环误差系统如图8.5.1所示，信号$\bm{\xi}_1$定义为
\begin{equation}\label{eq:8.5.12}
\bm{\xi}_1 = \tilde{\bm{W}}^T\hat{\sigma}'\hat{\bm{V}}^T\bm{x} + \tilde{\bm{W}}^T\hat{\sigma} + \bm{w} + \bm{\tau}_d
\end{equation}
注意神经网络在误差动力学中的作用，其中它被分解为两个有效块，出现在典型的反馈配置中。

我们现在揭示直接反向传播调整~\eqref{eq:8.5.7}产生的无源性性质。为了证明这个算法，使用与~\eqref{eq:8.5.2}不同的误差动力学形式，使得在图8.5.1中$\bm{\xi}_1 = \tilde{\bm{W}}^T\sigma(\bm{V}^T\bm{x}) + \bm{w} + \bm{\tau}_d$。

\textbf{定理8.5-3}：（反向传播神经网络调整算法的无源性）。

简单反向传播权重调整算法~\eqref{eq:8.5.7}使从$-\hat{\bm{W}}^T\sigma(\bm{V}^T\bm{x})$到$\hat{\sigma}'\hat{\bm{V}}^T\bm{x}$的映射和从$\bm{r}(t)$到$\hat{\bm{W}}^T\hat{\sigma}'\hat{\bm{V}}^T\bm{x}$的映射都是被动映射。

\textbf{证明}：

相对于$\bm{W}$、$\bm{V}$的动力学为
\begin{equation}
\dot{\tilde{\bm{W}}} = -\bm{F}\hat{\sigma}'\hat{\bm{V}}^T\bm{x}\bm{r}^T, \quad \dot{\tilde{\bm{V}}} = -\bm{G}\bm{x}(\hat{\bm{W}}^T\hat{\sigma}')^T\bm{r}
\end{equation}

1. 选择非负函数
\begin{equation}
L = \frac{1}{2}\text{tr}\{\tilde{\bm{W}}^T \bm{F}^{-1} \tilde{\bm{W}}\}
\end{equation}

并沿轨迹(1)评估$\dot{L}$得
\begin{equation}
\dot{L} = \text{tr}\{\tilde{\bm{W}}^T \bm{F}^{-1} \dot{\tilde{\bm{W}}}\} = -\text{tr}\{\tilde{\bm{W}}^T \hat{\sigma}'\hat{\bm{V}}^T\bm{x}\bm{r}^T\} = \langle -\hat{\bm{W}}^T\sigma(\bm{V}^T\bm{x}), \hat{\sigma}'\hat{\bm{V}}^T\bm{x} \rangle
\end{equation}

这是功率形式~\eqref{eq:8.5.8}。

2. 选择非负函数
\begin{equation}
L = \frac{1}{2}\text{tr}\{\tilde{\bm{V}}^T \bm{G}^{-1} \tilde{\bm{V}}\}
\end{equation}

并沿轨迹(2)评估得
\begin{equation}
\dot{L} = \text{tr}\{\tilde{\bm{V}}^T \bm{G}^{-1} \dot{\tilde{\bm{V}}}\} = -\text{tr}\{\tilde{\bm{V}}^T \bm{x}(\hat{\bm{W}}^T\hat{\sigma}')^T\bm{r}\} = \langle \bm{r}, \hat{\bm{W}}^T\hat{\sigma}'\hat{\bm{V}}^T\bm{x} \rangle
\end{equation}

这是功率形式。

因此，图8.5.1中的机器人误差系统是状态严格被动的（SSP），权重误差块是被动的；这保证了闭环系统的无源性。使用无源性定理~\cite{slotine1991}，现在可以得出，只要外部输入有界，每个块的输入/输出信号都是有界的。

不幸的是，虽然是被动的，但闭环系统不能保证是SSP，因此当扰动非零时，这不产生权重块内部状态（即$\bm{W}$、$\bm{V}$）的有界性，除非这些块是可观测的，即持续激励（PE）。

下一个结果显示为什么使用定理8.5.1中给出的修改权重更新算法不需要PE条件。

\textbf{定理8.5-4}：（增强反向传播神经网络调整算法的SSP）。

定理8.5.1中的修改权重调整算法使从$-\hat{\bm{W}}^T\sigma(\bm{V}^T\bm{x})$到$\hat{\sigma}'\hat{\bm{V}}^T\bm{x}$的映射和从$\bm{r}(t)$到$\hat{\bm{W}}^T\hat{\sigma}'\hat{\bm{V}}^T\bm{x}$的映射都是状态严格被动（SSP）映射。

\textbf{证明}：

使用定理8.5.1中的调整算法，相对于$\bm{W}$、$\bm{V}$的动力学由下式给出
\begin{align}
\dot{\tilde{\bm{W}}} &= -\bm{F}\hat{\sigma}'\hat{\bm{V}}^T\bm{x}\bm{r}^T + \bm{F}\hat{\sigma}'\tilde{\bm{V}}^T\bm{x}\bm{r}^T + \kappa\|\bm{r}\|\bm{F}\hat{\bm{W}} \\
\dot{\tilde{\bm{V}}} &= -\bm{G}\bm{x}(\hat{\bm{W}}^T\hat{\sigma}')^T\bm{r} + \kappa\|\bm{r}\|\bm{G}\hat{\bm{V}}
\end{align}

1. 选择非负函数
\begin{equation}
L = \frac{1}{2}\text{tr}\{\tilde{\bm{W}}^T \bm{F}^{-1} \tilde{\bm{W}}\}
\end{equation}

并评估$\dot{L}$得
\begin{align}
\dot{L} &= \text{tr}\{\tilde{\bm{W}}^T \bm{F}^{-1} \dot{\tilde{\bm{W}}}\} \\
&= \text{tr}\{\tilde{\bm{W}}^T(-\hat{\sigma}'\hat{\bm{V}}^T\bm{x}\bm{r}^T + \hat{\sigma}'\tilde{\bm{V}}^T\bm{x}\bm{r}^T + \kappa\|\bm{r}\|\hat{\bm{W}})\} \\
&= \langle -\hat{\bm{W}}^T\sigma(\bm{V}^T\bm{x}), \hat{\sigma}'\hat{\bm{V}}^T\bm{x} \rangle + \text{tr}\{\tilde{\bm{W}}^T\hat{\sigma}'\tilde{\bm{V}}^T\bm{x}\bm{r}^T\} + \kappa\|\bm{r}\|\text{tr}\{\tilde{\bm{W}}^T(\bm{W}-\tilde{\bm{W}})\}
\end{align}

由于
\begin{align}
\text{tr}\{\tilde{\bm{W}}^T\hat{\sigma}'\tilde{\bm{V}}^T\bm{x}\bm{r}^T\} &= \langle \bm{r}, \tilde{\bm{W}}^T\hat{\sigma}'\tilde{\bm{V}}^T\bm{x} \rangle = \langle \bm{r}, \hat{\bm{W}}^T\hat{\sigma}'\tilde{\bm{V}}^T\bm{x} \rangle - \langle \bm{r}, \bm{W}^T\hat{\sigma}'\tilde{\bm{V}}^T\bm{x} \rangle \\
&= \langle \bm{r}, \hat{\bm{W}}^T\hat{\sigma}'\tilde{\bm{V}}^T\bm{x} \rangle - \langle \bm{r}, \hat{\bm{W}}^T\hat{\sigma}'\tilde{\bm{V}}^T\bm{x} \rangle = 0
\end{align}

得到
\begin{equation}
\dot{L} = \langle -\hat{\bm{W}}^T\sigma(\bm{V}^T\bm{x}), \hat{\sigma}'\hat{\bm{V}}^T\bm{x} \rangle + \kappa\|\bm{r}\|\text{tr}\{\tilde{\bm{W}}^T\bm{W}\} - \kappa\|\bm{r}\|\|\tilde{\bm{W}}\|_F^2
\end{equation}

这是功率形式，最后一项在$\|\tilde{\bm{W}}\|_F$中是二次的。

2. 选择非负函数
\begin{equation}
L = \frac{1}{2}\text{tr}\{\tilde{\bm{V}}^T \bm{G}^{-1} \tilde{\bm{V}}\}
\end{equation}

并评估$\dot{L}$得
\begin{equation}
\dot{L} = \text{tr}\{\tilde{\bm{V}}^T \bm{G}^{-1} \dot{\tilde{\bm{V}}}\} = \langle \bm{r}, \hat{\bm{W}}^T\hat{\sigma}'\tilde{\bm{V}}^T\bm{x} \rangle + \kappa\|\bm{r}\|\text{tr}\{\tilde{\bm{V}}^T\bm{V}\} - \kappa\|\bm{r}\|\|\tilde{\bm{V}}\|_F^2
\end{equation}

这是功率形式，最后一项在$\|\tilde{\bm{V}}\|_F$中是二次的。

机器人动力学和权重调整块都是SSP，这保证了闭环系统的SSP，因此内部状态的范数以输送到每个块的功率为界。那么，输入/输出信号的有界性保证了在没有可观测性或PE要求的情况下状态有界性。

\textbf{被动神经网络和鲁棒神经网络的定义}。如果调整在误差形式中保证权重调整子系统的无源性，则我们将动态调整的神经网络定义为被动的。那么，需要额外的PE条件来保证权重的有界性。这个PE条件通常以神经网络隐层输出为条件。如果调整在误差形式中保证权重调整子系统的SSP，则我们将动态调整的神经网络定义为鲁棒的。那么，不需要额外的PE条件来保证权重的有界性。注意（1）此外还需要开环对象误差系统的SSP以保证跟踪稳定性，（2）神经网络无源性性质依赖于所使用的权重调整算法。

\textbf{1层神经网络控制器的无源性性质}

类似地，FLNN控制器调整算法使系统被动，因此需要额外的PE条件来验证神经网络权重的内部稳定性。另一方面，定理8.4.1中的增强调整算法产生SSP，因此不需要PE。

\section{小结}

为一般工业拉格朗日运动系统（以刚性机器人手臂为特征）给出了神经网络控制器设计算法。设计程序基于严格的非线性推导和稳定性证明，产生具有神经网络在某些环路中的多环智能控制结构。给出了不需要复杂初始化程序或任何离线学习阶段、实时在线工作、提供保证跟踪和有界神经网络权重和控制信号的神经网络权重调整算法。这里给出的神经网络控制器是无模型控制器，因为它们对规定类别的任何系统都有效，无需大量建模和初步分析来寻找"回归矩阵"。与标准机器人自适应控制器不同，它们不需要参数线性化（另见~\cite{colbaugh1994}，它不需要LIP）。适当的设计允许神经网络控制器避免诸如持续激励和确定性等价等要求。本文定义和研究了神经网络无源性和鲁棒性性质。

\begin{thebibliography}{99}

\bibitem{astrom1989} K.J.\AA str\"om and B.Wittenmark, \textit{Adaptive Control}, Addision Wesley, Reading, MA, 1989.

\bibitem{barron1993} Barron, A.R., ``Universal approximation bounds for superpositions of a sigmoidal function,'' \textit{IEEE Trans. Info. Theory}, vol. 39, no. 3, pp. 930--945, May 1993.

\bibitem{chen1992} F.-C.Chen and H.K.Khalil, ``Adaptive control of nonlinear systems using neural networks,'' \textit{Int. J. Control}, vol. 55, no. 6, pp. 1299--1317, 1992.

\bibitem{chen1994} F.-C.Chen and C.-C.Liu, ``Adaptively controlling nonlinear continuous-time systems using multilayer neural networks,'' \textit{IEEE Trans. Automat. Control}, vol. 39, no. 6, pp. 1306--1310, June 1994.

\bibitem{colbaugh1994} R.Colbaugh, H.Seraji, and K.Glass, ``A new class of adaptive controllers for robot trajectory tracking,'' \textit{J. Robotic Systems}, vol. 11, no. 8, pp. 761--772, 1994.

\bibitem{craig1988} Craig, J.J., \textit{Adaptive Control of Mechanical Manipulators}, Reading, MA: Addison-Wesley, 1988.

\bibitem{dawson1990} Dawson, D.M., Z.Qu, F.L.Lewis, and J.F.Dorsey, ``Robust control for the tracking of robot motion,'' \textit{Int. J. Control}, vol. 52, no. 3, pp. 581--595, 1990.

\bibitem{goodwin1984} Goodwin, C.G., and K.S.Sin, \textit{Adaptive Filtering, Prediction, and Control}, Prentice-Hall, New Jersey, 1984.

\bibitem{gorinevsky1995} D.Gorinevsky, ``On the persistency of excitation in radial basis function network identification of nonlinear systems,'' \textit{IEEE Trans. Neural Networks}, vol. 6, no. 5, pp. 1237--1244, Sept. 1995.

\bibitem{gutierrez1998} L.B.Gutierrez and F.L.Lewis, ``Implementation of a neural net tracking controller for a single flexible link: comparison with PD and PID controllers,'' \textit{IEEE Trans. Industrial Electronics}, to appear, April 1998.

\bibitem{haykin1994} S.Haykin, \textit{Neural Networks}, IEEE Press, New York, 1994.

\bibitem{hornik1989} K.Hornik, M.Stinchombe, and H.White, ``Multilayer feedforward networks are universal approximators,'' \textit{Neural Networks}, vol. 2, pp. 359--366, 1989.

\bibitem{ioannou1984} P.A.Ioannou and P.V.Kokotovic, ``Instability analysis and improvement of robustness of adaptive control,'' \textit{Automatica}, vol. 20, no. 5, pp. 583--594, 1984.

\bibitem{kim1998} Y.H.Kim and F.L.Lewis, \textit{High-Level Feedback Control with Neural Networks}, World Scientific, Singapore, 1998.

\bibitem{kreisselmeier1986} Kreisselmeier, G., and B.Anderson, ``Robust model reference adaptive control,'' \textit{IEEE Trans. Automat. Control}, vol. AC-31, no. 2, pp. 127--133, Feb. 1986.

\bibitem{landau1979} Y.D.Landau, \textit{Adaptive Control}, Marcel Dekker, Basel, 1979.

\bibitem{lewis1999b} F.L.Lewis, ``Nonlinear Network Structures for Feedback Control'', \textit{Asian Journal of Control}, vol. 1, no. 4, pp.205--228, December 1999.

\bibitem{lewis1993} F.L.Lewis, C.T.Abdallah, and D.M.Dawson, \textit{Control of Robot Manipulators}, New York: Macmillan, 1993.

\bibitem{lewis1999} F.L.Lewis, S.Jagannathan, and A.Yesildirek, \textit{Neural Network Control of Robot Manipulators and Nonlinear Systems}, Taylor and Francis, London, 1999.

\bibitem{lewis1998b} F.L.Lewis and Y.H.Kim, ``Neural Networks for Feedback Control,'' \textit{Encyclopedia of Electrical and Electronics Engineering}, ed. J.G.Webster, New York: Wiley, 1998. To appear.

\bibitem{lewis1993b} Lewis, F.L., K.Liu, and A.Yesildirek, ``Neural net robot controller with guaranteed tracking performance,'' \textit{Proc. Int. Symp. Intelligent Control}, pp. 225--231, Chicago., Aug. 1993.

\bibitem{lewis1995} Lewis, F.L., K.Liu, and A.Yesildirek, ``Neural net robot controller with guaranteed tracking performance,'' \textit{IEEE Trans. Neural Networks}, vol. 6, no. 3, pp. 703--715, 1995.

\bibitem{lewis1997} F.L.Lewis, K.Liu, and S.Commuri, ``Neural networks and fuzzy logic systems for robot control,'' in \textit{Fuzzy Logic and Neural Network Applications}, ed. F.Wang, World Scientific Pub., to appear, 1997.

\bibitem{lewis1996} F.L.Lewis, A.Yesildirek, and K.Liu, ``Multilayer neural net robot controller: structure and stability proofs,'' \textit{IEEE Trans. Neural Networks}, vol. 7, no. 2, pp. 1--12, Mar. 1996.

\bibitem{miller1991} W.T.Miller, R.S.Sutton, P.J.Werbos, ed., \textit{Neural Networks for Control}, Cambridge: MIT Press, 1991.

\bibitem{narendra1991} K.S.Narendra, ``Adaptive Control Using Neural Networks,'' Neural Networks for Control, pp 115--142. ed. W.T.Miller, R.S.Sutton, P.J.Werbos, Cambridge: MIT Press, 1991.

\bibitem{narendra1987} K.S.Narendra and A.M.Annaswamy, ``A new adaptive law for robust adaptive control without persistent excitation,'' \textit{IEEE Trans. Automat. Control}, vol. 32, pp. 134--145, Feb., 1987.

\bibitem{narendra1991b} Narendra, K.S., and K. Parthasarathy, ``Gradient methods for the optimization of dynamical systems containing neural networks,'' \textit{IEEE Trans. Neural Networks}, vol. 2, no. 2, pp. 252--262, Mar. 1991.

\bibitem{polycarpou1996} M.M.Polycarpou, ``Stable adaptive neural control scheme for nonlinear systems,'' \textit{IEEE Trans. Automat. Control}, vol. 41, no. 3, pp. 447--451, Mar. 1996.

\bibitem{polycarpou1991} M.M.Polycarpou and P.A.Ioannou, ``Identification and control using neural network models: design and stability analysis,'' Tech. Report 91--09--01, Dept. Elect. Eng. Sys., Univ. S. Cal., Sept. 1991.

\bibitem{rovithakis1994} G.A.Rovithakis and M.A. Christodoulou, ``Adaptive control of unknown plants using dynamical neural networks,'' \textit{IEEE Trans. Systems, Man, and Cybernetics}, vol. 24, no. 3, pp. 400--412, Mar. 1994.

\bibitem{sadegh1993} N.Sadegh, ``A perceptron network for functional identification and control of nonlinear systems,'' \textit{IEEE Trans. Neural Networks}, vol. 4, no. 6, pp. 982--988, Nov. 1993.

\bibitem{sanner1991} R.M.Sanner and J.-J.E.Slotine, ``Stable adaptive control and recursive identification using radial gaussian networks,'' \textit{Proc. IEEE Conf. Decision and Control}, Brighton, 1991.

\bibitem{seron1994} Seron, M.M., D.J.Hill, and A.L.Fradkov, ``Adaptive passification of nonlinear systems,'' \textit{Proc. IEEE Conf. Decision and Control}, pp. 190--195, Dec. 1994.

\bibitem{slotine1988} Slotine, J.-J.E., ``Putting physics in control the example of robotics,'' \textit{IEEE Control Systems Magazine}, pp. 12--17, Dec. 1988.

\bibitem{slotine1991} J.-J.Slotine and W.Li, \textit{Applied Nonlinear Control}, Prentice Hall, New Jersey, 1991.

\bibitem{spong1989} Spong, M.W., and M.Vidyasagar, \textit{Robot Dynamics and Control}, New York: Wiley, 1989.

\bibitem{wang1997} L.-X.Wang, \textit{A course in fuzzy systems and control}, Prentice-Hall, New Jersey, 1997.

\bibitem{werbos1989} P.J.Werbos, ``Back propagation: past and future,'' \textit{Proc. 1988 Int. Conf. Neural Nets}, vol. 1, pp. I343-I353, 1989.

\bibitem{werbos1992} P.J.Werbos, ``Neurocontrol and supervised learning: an overview and evaluation,'' in \textit{Handbook of Intelligent Control}, ed. D.A. White and D.A.Sofge, New York: Van Nostrand Reinhold, 1992.

\bibitem{white1992} D.A.White and D.A.Sofge, ed. \textit{Handbook of Intelligent Control}, New York: Van Nostrand Reinhold, 1992.

\end{thebibliography}

\ifstandalone
  \end{document}
\fi
